%! Inverse Trigonometric Functions — Unit 7, Lesson 7.2 — Exit Ticket (Answer Key)
\documentclass[10pt]{article}
\usepackage{precalculus-article}
\usepackage{precalculus-key}   % pulls in -boxes; do NOT also load -boxes
\newcommand{\TallMath}[1]{$\displaystyle #1\rule[-1.4em]{0pt}{3.2em}$}

\begin{document}

\pageheader{Unit 7, Lesson 7.2}{Exit Ticket --- Answer Key}

\vspace{0.12in}

\begin{notesbox}{Exit Ticket: Inverse Trigonometric Functions --- Key}

\textbf{1.} Evaluate each expression.  Give the exact answer in radians.

\renewcommand{\arraystretch}{2.0}
\begin{tabularx}{\linewidth}{@{} X X @{}}
  \textbf{(a)} $\arccos\!\left(-\dfrac{\sqrt{3}}{2}\right) = $
    \ans{$\dfrac{5\pi}{6}$}
  &
  \textbf{(b)} $\arctan\!\left(-\sqrt{3}\right) = $
    \ans{$-\dfrac{\pi}{3}$}
\end{tabularx}

\vspace{0.14in}

\textbf{2.} Explain why the domain of sine must be restricted.

\ans{The sine function is periodic and many-to-one: the same output value
is produced by infinitely many input angles.  Because it fails the horizontal
line test, it does not have an inverse unless we restrict its domain to an
interval on which it is one-to-one, such as $[-\tfrac{\pi}{2},\tfrac{\pi}{2}]$.}

\end{notesbox}

\end{document}
